% NCD-CIE v5-reviewed — Trimmed to 14 pages (LNCS format)
% Non-Communicable Disease Causal Inference Engine
\documentclass[runningheads]{llncs}
\usepackage[T1]{fontenc}
\usepackage{graphicx}
\usepackage{amsmath,amssymb}
\usepackage{booktabs}
\usepackage{algorithm}
\usepackage{algorithmic}
\usepackage{hyperref}
\usepackage{url}
\usepackage{multirow}
\usepackage{xcolor}

\titlerunning{NCD-CIE: Causal Inference Engine for NCD Risk}

\begin{document}

\title{NCD-CIE: A Causal Knowledge Graph and Counterfactual\\Intervention Engine for Non-Communicable Disease\\Risk Prediction}

\author{Anonymous Submission}
\authorrunning{Anonymous}
\institute{Anonymous Institution}

\maketitle

%--------------------------------------------------------------
\begin{abstract}
Non-communicable diseases (NCDs) such as cardiovascular disease (CVD),
type~2 diabetes mellitus (T2DM), and chronic kidney disease (CKD)
account for over 74\% of global mortality.  Existing risk calculators
(e.g., Framingham, QRISK3) rely on associational models that
cannot answer \emph{``what if''} questions about interventions.  We
present \textbf{NCD-CIE}, an open-source \emph{Causal Inference
Engine} that integrates an expert-curated causal knowledge graph (107
directed edges across 8 clinical domains), a logistic-link risk
prediction module with cross-source uncertainty quantification, and a
topological counterfactual intervention simulator grounded in Pearl's
structural causal model framework.  We formally define the causal
knowledge graph $\mathcal{G}=(V,E,W,\varepsilon)$ and position
NCD-CIE at rung~3 (counterfactuals) of Pearl's \emph{ladder of
causation}.  The system is validated on synthetic cohorts
($n=5{,}000$; C-statistics 0.76--0.82), NHANES 2017--2020
($n=8{,}291$; Pearson $r=0.91$ vs.\ Framingham, calibration slope
0.94), and face-validity comparisons against four landmark RCTs.  We
further present data-driven graph validation via the PC algorithm,
subgroup fairness analysis, and decision curve analysis demonstrating
positive net clinical benefit across thresholds of 5--30\%.
Population transportability to Thai/Asian cohorts is discussed with
reference to Bareinboim--Pearl transportability theory.

\keywords{Causal inference \and Knowledge graph \and
Non-communicable disease \and Risk prediction \and
Counterfactual reasoning \and Clinical decision support}
\end{abstract}

%==============================================================
\section{Introduction}\label{sec:intro}

Non-communicable diseases---cardiovascular disease (CVD), type~2
diabetes mellitus (T2DM), chronic kidney disease (CKD), and their
comorbid clusters---impose an enormous global burden, responsible for
41 million deaths annually~\cite{who2021ncd}.  Despite well-validated
risk scores~\cite{dagostino2008framingham, hippisley2017qrisk3,
score22021}, current tools share a fundamental limitation: they
estimate \emph{associational} risk but cannot reason about
\emph{causal} mechanisms or simulate the effect of hypothetical
interventions.

A clinician who asks ``What would this patient's 10-year CVD risk be
if they lowered their LDL-C by 1\,mmol/L and lost 5\,kg?'' receives
no answer from Framingham or QRISK3.  Answering such questions
requires moving beyond rung~1 (association) of Pearl's ladder of
causation~\cite{pearl2018book} to rung~3 (counterfactuals).

We propose \textbf{NCD-CIE} (Non-Communicable Disease Causal Inference
Engine), combining three integrated components:
\begin{enumerate}
\item A \textbf{causal knowledge graph} $\mathcal{G}=(V,E,W,\varepsilon)$
  encoding directed causal relationships among biomarkers, lifestyle
  factors, and disease endpoints across 8 clinical domains, curated
  according to Bradford Hill criteria~\cite{hill1965environment} with
  inter-rater reliability $\kappa=0.78$.
\item A \textbf{risk prediction engine} using logistic-link scoring with
  cross-source uncertainty quantification.
\item A \textbf{what-if intervention simulator} that propagates
  counterfactual changes through the causal graph using topological
  sorting, producing clinically interpretable risk deltas.
\end{enumerate}

\noindent\textbf{Contributions.}
(C1)~A formally defined causal knowledge graph for NCD risk with 107
expert-curated edges, positioned at rung~3 of the ladder of causation.
(C2)~A logistic-link risk engine with cross-source uncertainty fusion.
(C3)~A topological counterfactual simulator for personalised what-if queries.
(C4)~Multi-source validation: synthetic cohorts, NHANES external validation,
and face-validity alignment with four landmark RCTs.
(C5)~Data-driven graph validation via the PC algorithm and subgroup fairness
analysis.
(C6)~A transportability analysis for Thai/Asian populations.

%==============================================================
\section{Related Work}\label{sec:related}

\paragraph{NCD Risk Prediction.}
Traditional calculators estimate 10-year event probabilities using Cox
proportional hazards.  The Framingham Risk
Score~\cite{dagostino2008framingham} pioneered this approach;
QRISK3~\cite{hippisley2017qrisk3} extended predictors to include
ethnicity and comorbidities; SCORE2~\cite{score22021} recalibrated
for European regions.  While achieving C-statistics of 0.71--0.78,
none support causal or counterfactual reasoning.  Machine learning
approaches have improved discrimination
modestly~\cite{alaa2019attentive} but remain fundamentally
associational.

\paragraph{Causal Inference Frameworks.}
Pearl's SCM framework~\cite{pearl2009causality} provides the
theoretical foundation for interventional and counterfactual reasoning
via the \emph{do}-calculus.  Software implementations include
\textbf{DoWhy}~\cite{sharma2020dowhy}, targeting observational study
analysis rather than real-time clinical decision support, and
\textbf{CausalNex}~\cite{beaumont2021causalnex}, which combines
Bayesian networks with domain knowledge but lacks health-specific
ontology and intervention simulation.

\paragraph{Health Knowledge Graphs.}
\textbf{PrimeKG}~\cite{chandak2023primekg} integrates 20 biomedical
resources with over 100{,}000 nodes but encodes primarily
associational relationships.  Barab\'{a}si's network medicine
framework~\cite{barabasi2011network} established the theoretical
basis for disease--gene--drug networks but does not address
individual-level risk prediction.  NCD-CIE differs in three respects:
(i)~edges are explicitly \emph{causal} and \emph{directed}, curated
against Bradford Hill criteria; (ii)~each edge carries a quantitative
effect size with uncertainty; and (iii)~the graph is purpose-built for
individual-level counterfactual intervention simulation.

%==============================================================
\section{System Architecture}\label{sec:arch}

NCD-CIE follows a modular, layered architecture (Fig.~\ref{fig:arch})
comprising four tiers: data ingestion, causal knowledge graph,
inference engine, and presentation layer.

\begin{figure}[t]
\centering
\fbox{\parbox{0.9\textwidth}{\centering
\textbf{[Figure 1: System Architecture Diagram]}\\[4pt]
\small Data Sources $\rightarrow$ Ingestion Layer $\rightarrow$
Causal KG $\rightarrow$ Risk Engine + What-If Simulator
$\rightarrow$ REST API / Dashboard}}
\caption{NCD-CIE system architecture.  Data flows from heterogeneous
sources through the causal knowledge graph to the risk prediction and
intervention simulation modules.}\label{fig:arch}
\end{figure}

\subsection{Formal Definition of the Causal Knowledge Graph}
\label{sec:formal}

We define the NCD-CIE causal knowledge graph as a weighted directed
acyclic graph (DAG):
\begin{equation}\label{eq:graph}
\mathcal{G} = (V, E, W, \varepsilon)
\end{equation}
where $V=\{v_1,\ldots,v_n\}$ is the set of $n$ biomarker and clinical
variable nodes, $E \subseteq V \times V$ is the set of directed causal
edges, $W: E \rightarrow \mathbb{R}$ assigns a causal effect weight
(log-odds ratio or standardised mean difference) to each edge, and
$\varepsilon: E \rightarrow \mathbb{R}^+$ assigns an uncertainty
(standard error) to each weight.

\paragraph{Positioning within Causal Frameworks.}
Table~\ref{tab:comparison} contrasts $\mathcal{G}$ with related
formalisms.  Unlike traditional knowledge graphs, $\mathcal{G}$
enforces strict causal directionality.  Unlike Bayesian networks,
$\mathcal{G}$ uses sparse parametric edge weights with uncertainty,
enabling efficient counterfactual propagation without full joint
distribution specification.  Unlike Pearl's full SCM, $\mathcal{G}$
operates with a lightweight linear-on-log-odds approximation
sufficient for clinical risk scoring.

\begin{table}[t]
\centering
\caption{Comparison of NCD-CIE's causal KG with related formalisms.}
\label{tab:comparison}
\small
\begin{tabular}{@{}lccccc@{}}
\toprule
\textbf{Property} & \textbf{Trad.\ KG} & \textbf{BN} & \textbf{SCM} & \textbf{DoWhy} & \textbf{NCD-CIE} \\
\midrule
Directed edges       & Partial & Yes & Yes & Yes & Yes \\
Causal semantics     & No      & Partial & Yes & Yes & Yes \\
Edge weights         & No      & CPTs & Equations & Estimated & $W \pm \varepsilon$ \\
Counterfactuals      & No      & No  & Yes & Limited & Yes \\
Real-time scoring    & No      & Slow & Slow & No & Yes \\
NCD-specific         & No      & No  & No  & No & Yes \\
Ladder rung          & 1       & 1--2 & 3   & 2  & 3 \\
\bottomrule
\end{tabular}
\end{table}

NCD-CIE operates at \textbf{rung~3} of Pearl's ladder of
causation~\cite{pearl2018book}: given an individual's observed profile
$(X=x, Y=y)$, the system answers counterfactual queries of the form
$P(Y_{x'}|X=x, Y=y)$.

\subsection{Clinical Domains}

The knowledge graph spans 8 clinical domains (Table~\ref{tab:domains}),
covering the major modifiable and non-modifiable risk factors for CVD,
T2DM, and CKD.

\begin{table}[t]
\centering
\caption{Eight clinical domains in the NCD-CIE knowledge graph.}
\label{tab:domains}
\small
\begin{tabular}{@{}clcc@{}}
\toprule
\textbf{\#} & \textbf{Domain} & \textbf{Nodes} & \textbf{Edges} \\
\midrule
1 & Lipid metabolism        & 8  & 14 \\
2 & Glycaemic regulation    & 7  & 12 \\
3 & Blood pressure          & 6  & 11 \\
4 & Renal function          & 5  & 9  \\
5 & Inflammatory markers    & 6  & 13 \\
6 & Anthropometric / body composition & 5 & 10 \\
7 & Lifestyle / behavioural & 8  & 18 \\
8 & Disease endpoints       & 6  & 20 \\
\midrule
  & \textbf{Total}          & \textbf{51} & \textbf{107} \\
\bottomrule
\end{tabular}
\end{table}

%==============================================================
\section{Causal Knowledge Graph}\label{sec:kg}

\subsection{Biomarker Ontology}\label{sec:ontology}

Each node $v \in V$ represents a measurable clinical quantity or
categorical risk factor.  Continuous biomarkers (e.g., LDL-C, HbA1c,
eGFR) are standardised to population $z$-scores using age- and
sex-specific reference distributions from NHANES~\cite{nhanes2020}.
Categorical variables are encoded as binary indicators.  Disease
endpoints are modelled as latent binary nodes with logistic-link
activation.  The ontology aligns with LOINC codes and maps to
SNOMED-CT concepts for EHR interoperability.

\subsection{Edge Construction and Bradford Hill Criteria}
\label{sec:edges}

Edges represent direct causal effects, curated from systematic
reviews, meta-analyses, Mendelian randomisation studies, and landmark
RCTs.  Each candidate edge undergoes a structured assessment against
Bradford Hill's nine criteria for causal
inference~\cite{hill1965environment}: strength, consistency,
specificity, temporality, biological gradient, plausibility,
coherence, experiment, and analogy.

Each edge is assigned an evidence grade $g \in \{A, B, C\}$:
\textbf{Grade~A} (high confidence) satisfies criteria 1--8, sourced
from RCTs or Mendelian randomisation (e.g., LDL-C $\rightarrow$ CVD);
\textbf{Grade~B} (moderate) satisfies criteria 1--7, supported by
large observational cohorts (e.g., CRP $\rightarrow$ CVD);
\textbf{Grade~C} (emerging) satisfies criteria 1, 2, 4, 6, with
weight down-scaled by 0.5.  Two independent clinical reviewers
assessed each edge ($\kappa = 0.78$, substantial agreement);
disagreements were resolved by a third reviewer.  The intraclass
correlation coefficient for edge weights was ICC(2,1)\,=\,0.84.

\subsection{Cycle Handling}\label{sec:cycles}

Biological feedback loops (e.g., insulin resistance
$\leftrightarrow$ obesity) are handled by applying Tarjan's algorithm
to identify strongly connected components (SCCs), collapsing each into
a super-node with internal dynamics modelled by fixed-point iteration
(convergence threshold $10^{-4}$, max 50 iterations; all current SCCs
converge within 8).  Ablation (Section~\ref{sec:ablation}) confirms
SCC handling has negligible impact on aggregate concordance.

%==============================================================
\section{Risk Prediction Engine}\label{sec:risk}

\subsection{Logistic-Link Scoring}\label{sec:logistic}

For each disease endpoint $d \in \{\text{CVD}, \text{T2DM}, \text{CKD}\}$,
the baseline 10-year risk is:
\begin{equation}\label{eq:risk}
R_d = \sigma\!\left(\beta_{0,d} + \sum_{(v_i, d) \in E} W_{(v_i,d)} \cdot z_i\right)
\end{equation}
where $\sigma(\cdot)$ is the logistic sigmoid, $\beta_{0,d}$ is the
disease-specific intercept calibrated to population base rates, and
$z_i$ is the standardised biomarker value.  The logistic link was
chosen because it (i)~directly yields probability estimates, (ii)~composes
naturally with counterfactual propagation, and (iii)~for 10-year horizons
with event rates below 30\%, yields near-identical predictions to Cox
models~\cite{dagostino2008framingham}.

\subsection{Cross-Source Uncertainty Quantification}\label{sec:uncert}

Edge weights are modelled as $W_e \sim \mathcal{N}(\hat{W}_e,
\varepsilon_e^2)$.  For edges with multiple source estimates, we apply
inverse-variance weighted meta-analysis:
\begin{equation}\label{eq:meta}
\hat{W}_e = \frac{\sum_k \hat{W}_e^{(k)} / (\varepsilon_e^{(k)})^2}
{\sum_k 1 / (\varepsilon_e^{(k)})^2}, \quad
\varepsilon_e^2 = \frac{1}{\sum_k 1 / (\varepsilon_e^{(k)})^2}
\end{equation}
Risk uncertainty is propagated via first-order Taylor expansion:
\begin{equation}\label{eq:riskvar}
\mathrm{Var}(R_d) \approx \sum_{(v_i,d) \in E}
\left(\frac{\partial R_d}{\partial W_{(v_i,d)}}\right)^2 \varepsilon_{(v_i,d)}^2
\end{equation}
yielding 95\% credible intervals $R_d \pm 1.96\sqrt{\mathrm{Var}(R_d)}$.

\subsection{Composite NCD Risk Score}\label{sec:composite}

A composite score integrating all endpoints is computed as:
\begin{equation}\label{eq:composite}
R_{\mathrm{NCD}} = 1 - \prod_{d} (1 - R_d)
\end{equation}
under a conditional independence assumption, yielding the probability
of at least one NCD event within the prediction horizon.

%==============================================================
\section{What-If Intervention Simulator}\label{sec:whatif}

Given an individual's observed profile $\mathbf{x}$ and intervention
$\mathrm{do}(v_j = x_j')$, the counterfactual risk is
$R_d^{\mathrm{CF}} = R_d(\mathbf{x}^{\mathrm{CF}})$, where
$\mathbf{x}^{\mathrm{CF}}$ is obtained by propagating
$\Delta_j = x_j' - x_j$ through downstream nodes.  This corresponds
to the effect of treatment on the treated
(ETT)~\cite{pearl2009causality}:
$\mathrm{ETT} = E[Y_{x'} - Y_x \mid X=x, Y=y]$.

\subsection{Topological Cascade Algorithm}\label{sec:cascade}

The counterfactual profile is computed by
Algorithm~\ref{alg:cascade}, processing nodes in topological order.

\begin{algorithm}[t]
\caption{Topological Counterfactual Cascade}
\label{alg:cascade}
\begin{algorithmic}[1]
\REQUIRE Graph $\mathcal{G}$, profile $\mathbf{x}$, intervention
  $\mathrm{do}(v_j = x_j')$, depth $d_{\max}$, attenuation $\gamma$
\ENSURE Counterfactual profile $\mathbf{x}^{\mathrm{CF}}$
\STATE $\mathbf{x}^{\mathrm{CF}} \leftarrow \mathbf{x}$;
  $x_j^{\mathrm{CF}} \leftarrow x_j'$;
  $\Delta_j \leftarrow x_j' - x_j$
\STATE $\mathcal{S} \leftarrow \text{TopologicalSort}(\mathcal{G})$
\FOR{each $v_k$ in $\mathcal{S}$ after $v_j$}
  \STATE $\delta_k \leftarrow \sum_{(v_p, v_k) \in E} W_{(v_p,v_k)} \cdot (x_p^{\mathrm{CF}} - x_p) \cdot \gamma^{\text{depth}(v_j,v_p)}$ for depth $< d_{\max}$
  \STATE $x_k^{\mathrm{CF}} \leftarrow x_k + \delta_k$
\ENDFOR
\RETURN $\mathbf{x}^{\mathrm{CF}}$
\end{algorithmic}
\end{algorithm}

The attenuation factor $\gamma \in (0,1)$ controls decay of indirect
effects with graph distance.  Default values ($\gamma = 0.7$,
$d_{\max} = 3$) were selected via sensitivity analysis
(Section~\ref{sec:sensitivity}).

\paragraph{Lifestyle Intervention Mapping.}
To bridge clinical biomarkers and actionable changes, NCD-CIE maps
everyday interventions to $\mathrm{do}(\cdot)$ operations: physical
activity (steps/day) $\rightarrow$ $\Delta$HDL-C, $\Delta$SBP,
$\Delta$BMI~\cite{naci2018comparative}; dietary changes $\rightarrow$
$\Delta$LDL-C, $\Delta$CRP; weight loss $\rightarrow$ cascading
effects on metabolic nodes; pharmacological interventions (statins,
SGLT2i, antihypertensives) $\rightarrow$ target-specific $\Delta$
values calibrated from meta-analyses.

%==============================================================
\section{Validation}\label{sec:validation}

\subsection{Synthetic Cohort Validation}\label{sec:synthetic}

A synthetic cohort ($n = 5{,}000$) was generated using the causal
graph's structural equations with Gaussian noise (70/30 train/test
split).  Results (Table~\ref{tab:synthetic}) show C-statistics of
0.76--0.82 and calibration slopes of 0.89--0.93.

\begin{table}[t]
\centering
\caption{Synthetic cohort validation ($n = 5{,}000$, 30\% test set).}
\label{tab:synthetic}
\small
\begin{tabular}{@{}lccc@{}}
\toprule
\textbf{Metric} & \textbf{CVD} & \textbf{T2DM} & \textbf{CKD} \\
\midrule
C-statistic        & 0.82 & 0.79 & 0.76 \\
Calibration slope  & 0.93 & 0.91 & 0.89 \\
Brier score        & 0.142 & 0.168 & 0.098 \\
\bottomrule
\end{tabular}
\end{table}

\subsection{NHANES External Validation}\label{sec:nhanes}

NCD-CIE was validated against the Framingham Risk Score using NHANES
2017--2020 data~\cite{nhanes2020} ($n = 8{,}291$ adults, age 30--75,
complete lab panel, no prior CVD).  Table~\ref{tab:nhanes} summarises results.

\begin{table}[t]
\centering
\caption{NHANES external validation: NCD-CIE vs.\ Framingham
($n = 8{,}291$).}
\label{tab:nhanes}
\small
\begin{tabular}{@{}lc@{}}
\toprule
\textbf{Metric} & \textbf{Value} \\
\midrule
Pearson $r$ (NCD-CIE vs.\ Framingham) & 0.91 \\
Calibration slope                       & 0.94 \\
Mean absolute difference                & 2.3 pp \\
Agreement within $\pm 5$\% absolute     & 87.4\% \\
\bottomrule
\end{tabular}
\end{table}

\subsection{Face Validity Against Landmark RCTs}\label{sec:faceval}

We simulated each trial's intervention on an NHANES-derived population
matching trial eligibility criteria (Table~\ref{tab:rct}).  All
simulated effects fall within the 95\% CIs of the corresponding RCT
results.  The DPP comparison shows the largest discrepancy
($-$11.3\% vs.\ $-$16.0\%), attributable to the DPP trial's intensive
lifestyle intervention affecting pathways beyond weight loss alone.

\begin{table}[t]
\centering
\caption{Face validity: NCD-CIE simulated vs.\ RCT observed effects.}
\label{tab:rct}
\small
\begin{tabular}{@{}llcc@{}}
\toprule
\textbf{Intervention} & \textbf{RCT} &
\textbf{NCD-CIE} & \textbf{RCT} \\
\midrule
Statin (LDL-C $-$1 mmol/L) & CTT~\cite{ctt2010}
  & $-$5.8\% & $-$5.4\% \\
Weight loss ($-$7\%) & DPP~\cite{knowler2002dpp}
  & $-$11.3\% & $-$16.0\% \\
SGLT2 inhibitor & CREDENCE~\cite{perkovic2019credence}
  & $-$3.1\% & $-$2.5\% \\
SBP $-$15 mmHg & SPRINT~\cite{sprint2015}
  & $-$4.2\% & $-$4.1\% \\
\bottomrule
\end{tabular}
\end{table}

\subsection{Sensitivity Analysis and Ablation}\label{sec:sensitivity}\label{sec:ablation}

Table~\ref{tab:sensitivity} shows the statin intervention effect is
moderately sensitive to $\gamma$ (range: $-$3.9\% to $-$8.1\%) but
robust to $d_{\max}$ beyond 3.  The default $\gamma = 0.7$ yields
the closest match to RCT evidence.

\begin{table}[t]
\centering
\caption{Sensitivity of $\Delta R_{\text{CVD}}$ (statin) to
hyperparameters, and ablation study (Pearson $r$ vs.\ Framingham).}
\label{tab:sensitivity}
\small
\begin{tabular}{@{}lcc|lc@{}}
\toprule
\textbf{Param.} & \textbf{Value} & \textbf{$\Delta R_{\text{CVD}}$} &
\textbf{Configuration} & \textbf{$r$} \\
\midrule
$\gamma$ & 0.5 & $-$3.9\% & Full model & 0.91 \\
$\gamma$ & 0.7 & $-$5.8\% & Grade A only & 0.89 \\
$\gamma$ & 0.9 & $-$8.1\% & No SCC handling & 0.91 \\
$d_{\max}$ & 2 & $-$5.1\% & Shuffled weights & 0.72$\pm$0.08 \\
$d_{\max}$ & 3 & $-$5.8\% & & \\
\bottomrule
\end{tabular}
\end{table}

The shuffled-weights control ($r = 0.72 \pm 0.08$) confirms that the
specific causal structure, not merely graph topology, drives
performance.

\subsection{Data-Driven Graph Validation}\label{sec:pcalg}

To assess whether the expert-curated structure is consistent with
data-driven causal discovery, we applied the PC
algorithm~\cite{spirtes2000causation} to NHANES using the same 51
variables.  The PC algorithm estimates a completed partially directed
acyclic graph (CPDAG) under the faithfulness assumption.

\begin{table}[t]
\centering
\caption{Expert-curated KG vs.\ PC-algorithm graph (NHANES).}
\label{tab:pcalg}
\small
\begin{tabular}{@{}lc@{}}
\toprule
\textbf{Metric} & \textbf{Value} \\
\midrule
Expert KG edges                     & 107 \\
PC-algorithm edges (CPDAG)          & 143 \\
Edge agreement (shared edges)       & 69 (64.5\%) \\
Expert-only edges                   & 38 \\
PC-only edges                       & 74 \\
Structural Hamming distance (SHD)   & 112 \\
\bottomrule
\end{tabular}
\end{table}

The 64.5\% edge agreement indicates that the majority of
expert-curated edges are data-corroborated.  The expert graph is more
conservative (107 vs.\ 143 edges), consistent with stricter Bradford
Hill criteria.  Expert-only edges largely represent well-established
pathways with small effect sizes below the PC algorithm's detection
threshold (e.g., moderate alcohol $\rightarrow$ HDL-C, uric acid
$\rightarrow$ CKD).  PC-only edges include plausible associations not
meeting Bradford Hill criteria (e.g., age $\rightarrow$ multiple
biomarkers, reflecting confounding rather than direct causation).

\subsection{Subgroup Fairness Analysis}\label{sec:subgroup}

Table~\ref{tab:subgroup} stratifies performance by demographics.
The maximum C-statistic gap across racial/ethnic groups is 0.05
(0.78--0.83), comparable to disparities reported for
Framingham and QRISK3~\cite{hippisley2017qrisk3}.

\begin{table}[t]
\centering
\caption{Subgroup analysis: C-statistics and concordance by
demographic stratum (NHANES, CVD).}
\label{tab:subgroup}
\small
\begin{tabular}{@{}llcc@{}}
\toprule
\textbf{Stratum} & \textbf{Subgroup} & \textbf{C-stat} & \textbf{$r$} \\
\midrule
\multirow{4}{*}{Race/Eth.}
  & NH White  & 0.83 & 0.93 \\
  & NH Black  & 0.79 & 0.88 \\
  & Hispanic  & 0.80 & 0.86 \\
  & Asian     & 0.78 & 0.83 \\
\midrule
\multirow{4}{*}{Age}
  & 30--44    & 0.77 & 0.86 \\
  & 45--54    & 0.81 & 0.90 \\
  & 55--64    & 0.84 & 0.94 \\
  & 65--75    & 0.82 & 0.92 \\
\midrule
\multirow{2}{*}{Sex}
  & Male      & 0.83 & 0.92 \\
  & Female    & 0.80 & 0.88 \\
\bottomrule
\end{tabular}
\end{table}

\subsection{Decision Curve Analysis}\label{sec:dca}

Decision curve analysis~\cite{vickers2006dca} compares NCD-CIE,
Framingham, and treat-all/treat-none strategies
(Table~\ref{tab:dca}).  NCD-CIE achieves higher net benefit at all
thresholds, with advantage increasing at higher thresholds.  The
net reclassification improvement is $\text{NRI}_{\text{overall}} =
+0.100$ (95\% CI: 0.041--0.159).

\begin{table}[t]
\centering
\caption{Decision curve analysis: net benefit ($\times 10^{-2}$) at
selected thresholds (NHANES, CVD).}
\label{tab:dca}
\small
\begin{tabular}{@{}lccc@{}}
\toprule
\textbf{Threshold} & \textbf{Treat All} & \textbf{Framingham} &
\textbf{NCD-CIE} \\
\midrule
5\%  & 4.2 & 5.1 & 5.3 \\
10\% & 2.8 & 4.3 & 4.6 \\
15\% & 1.1 & 3.5 & 3.9 \\
20\% & $-$0.3 & 2.8 & 3.2 \\
25\% & $-$1.5 & 2.1 & 2.6 \\
30\% & $-$2.4 & 1.5 & 2.0 \\
\bottomrule
\end{tabular}
\end{table}

%==============================================================
\section{Implementation and Transportability}\label{sec:impl}

NCD-CIE is implemented as a modular Python library (Python $\geq$
3.9) with components for graph management (\texttt{ncdcie.graph}),
risk prediction (\texttt{ncdcie.risk}), counterfactual simulation
(\texttt{ncdcie.whatif}), and a RESTful API (\texttt{ncdcie.api}).

\paragraph{Use Case.}
A 55-year-old male with LDL-C 4.2\,mmol/L, SBP 148\,mmHg, HbA1c
6.8\%, BMI 31, eGFR 68 receives: $R_{\text{CVD}} = 18.3\%$ (95\% CI:
14.1--22.5\%), $R_{\text{T2DM}} = 24.7\%$, $R_{\text{CKD}} = 12.1\%$.
Querying ``start a statin, lose 5\,kg, reduce sodium'' yields
$R_{\text{CVD}}^{\text{CF}} = 11.7\%$ ($-$6.6 pp, 36\% relative
reduction).

\subsection{Thai/Asian Transportability}\label{sec:thai}

Thai and Asian populations exhibit distinct NCD risk profiles: lower
BMI but higher visceral adiposity, different lipid distributions, and
earlier T2DM onset~\cite{aekplakorn2011thai}.  Following
Bareinboim--Pearl transportability
theory~\cite{bareinboim2016causal}, we assume structural invariance
of the graph topology while recalibrating edge weights and intercepts
using local cohort data.  Candidate sources include the EGAT
study~\cite{sritara2003egat} and Thai NHES~\cite{aekplakorn2011thai}.
Preliminary analysis suggests recalibrating the BMI$\rightarrow$T2DM
weight from 0.68 to 0.85 and adjusting the SBP$\rightarrow$CVD
intercept by $+$0.12 log-odds aligns predictions with EGAT observed
rates, demonstrating cross-population adaptability without structural
modification.

%==============================================================
\section{Safety and Ethical Considerations}\label{sec:ethics}

NCD-CIE is a \emph{decision-support} tool, not a diagnostic system.
Key safeguards include: (1)~all risk estimates include 95\% credible
intervals, with uncertainty prominently displayed; (2)~clinical
guardrails flag interventions producing physiologically implausible
values (e.g., LDL-C $< 0.5$\,mmol/L); (3)~every output includes a
disclaimer that results should not replace clinical judgement;
(4)~subgroup performance monitoring (Section~\ref{sec:subgroup})
triggers alerts when disparities exceed thresholds; (5)~the system
operates on de-identified data with no persistent storage, supporting
on-premise deployment for compliance with data protection regulations
(e.g., Thailand's PDPA); and (6)~the causal graph is fully
inspectable, enabling users to trace any estimate to its contributing
edges and evidence sources.

%==============================================================
\section{Discussion}\label{sec:discussion}

NCD-CIE demonstrates that a causal knowledge graph with logistic-link
scoring and topological counterfactual propagation can produce
clinically plausible predictions that closely align with established
models ($r = 0.91$), simulate intervention effects consistent with
RCTs, and provide mechanistic transparency.

\paragraph{Strengths.}
The formal grounding in Pearl's SCM framework positions NCD-CIE at
rung~3 of the ladder of causation, enabling queries that no existing
risk calculator can answer.  The Bradford Hill--based curation
protocol provides a principled, reproducible methodology for edge
construction.  Data-driven validation via the PC algorithm
(Section~\ref{sec:pcalg}) provides independent structural
confirmation, and the subgroup analysis reveals performance
disparities for targeted improvement.

\paragraph{Limitations.}
(1)~The linear-on-log-odds assumption does not capture non-linearities
or interactions (e.g., age $\times$ sex $\times$ cholesterol); future
work could incorporate non-linear structural equations.
(2)~Expert curation is slow, subjective, and potentially incomplete.
While PC algorithm validation provides a data-driven check, the system
cannot discover novel causal pathways automatically.
(3)~The current graph is static; as new evidence accumulates, edges
must be manually updated.  Bayesian updating from streaming evidence
is a planned enhancement.
(4)~The composite score (Eq.~\ref{eq:composite}) assumes conditional
independence of endpoints given risk factors, which may underestimate
risk for correlated disease pathways (e.g., diabetic nephropathy).
(5)~NHANES validation uses cross-sectional data with imputed 10-year
risk rather than true longitudinal outcomes; prospective validation is
needed.
(6)~Performance varies by demographic group
(Section~\ref{sec:subgroup}); the recalibration framework addresses
this in principle but lacks prospective validation.

\paragraph{Future Work.}
Key directions include prospective validation in the EGAT cohort,
EHR integration via FHIR APIs, extension to cancer endpoints,
incorporation of polygenic risk scores, and a patient-facing mobile
application.

%==============================================================
\section{Conclusion}\label{sec:conclusion}

We presented NCD-CIE, a causal inference engine for NCD risk
prediction and counterfactual intervention simulation.  Grounded in a
formally defined causal knowledge graph
$\mathcal{G} = (V, E, W, \varepsilon)$ curated against Bradford Hill
criteria, NCD-CIE operates at rung~3 of Pearl's ladder of causation.
Validation across synthetic cohorts, NHANES, and four landmark RCTs
demonstrates predictive accuracy ($r = 0.91$, C-statistics
0.76--0.82), calibration quality (slopes 0.89--0.94), and clinical
face validity.  Data-driven graph validation confirms 64.5\% edge
agreement with the PC algorithm, subgroup analysis reveals acceptable
equity, and decision curve analysis shows positive net benefit across
5--30\% thresholds.  NCD-CIE is released as an open-source Python
library with a transportable architecture supporting
population-specific recalibration.

%==============================================================
\begin{thebibliography}{30}

\bibitem{who2021ncd}
World Health Organization:
Non-communicable diseases fact sheet (2021).
\url{https://www.who.int/news-room/fact-sheets/detail/noncommunicable-diseases}

\bibitem{dagostino2008framingham}
D'Agostino, R.B., Vasan, R.S., Pencina, M.J., et al.:
General cardiovascular risk profile for use in primary care.
Circulation \textbf{117}(6), 743--753 (2008)

\bibitem{hippisley2017qrisk3}
Hippisley-Cox, J., Coupland, C., Brindle, P.:
Development and validation of QRISK3.
BMJ \textbf{357}, j2099 (2017)

\bibitem{score22021}
SCORE2 Working Group:
SCORE2 risk prediction algorithms.
European Heart Journal \textbf{42}(25), 2439--2454 (2021)

\bibitem{pearl2009causality}
Pearl, J.:
Causality: Models, Reasoning, and Inference, 2nd edn.
Cambridge University Press (2009)

\bibitem{pearl2018book}
Pearl, J., Mackenzie, D.:
The Book of Why: The New Science of Cause and Effect.
Basic Books (2018)

\bibitem{hill1965environment}
Hill, A.B.:
The environment and disease: association or causation?
Proc.\ Royal Society of Medicine \textbf{58}(5), 295--300 (1965)

\bibitem{sharma2020dowhy}
Sharma, A., Kiciman, E.:
DoWhy: An end-to-end library for causal inference.
arXiv:2011.04216 (2020)

\bibitem{beaumont2021causalnex}
Beaumont, P., et al.:
CausalNex: Toolkit for causal reasoning with Bayesian networks (2021)

\bibitem{bareinboim2016causal}
Bareinboim, E., Pearl, J.:
Causal inference and the data-fusion problem.
PNAS \textbf{113}(27), 7345--7352 (2016)

\bibitem{chandak2023primekg}
Chandak, P., Huang, K., Zitnik, M.:
Building a knowledge graph to enable precision medicine.
Scientific Data \textbf{10}, 67 (2023)

\bibitem{barabasi2011network}
Barab\'{a}si, A.L., Gulbahce, N., Loscalzo, J.:
Network medicine: a network-based approach to human disease.
Nature Reviews Genetics \textbf{12}(1), 56--68 (2011)

\bibitem{alaa2019attentive}
Alaa, A.M., van der Schaar, M.:
Attentive state-space modeling of disease progression.
NeurIPS (2019)

\bibitem{nhanes2020}
Centers for Disease Control and Prevention:
NHANES 2017--2020.
\url{https://www.cdc.gov/nchs/nhanes/}

\bibitem{ctt2010}
CTT Collaboration:
Efficacy and safety of more intensive lowering of LDL cholesterol.
The Lancet \textbf{376}(9753), 1670--1681 (2010)

\bibitem{knowler2002dpp}
Knowler, W.C., et al.:
Reduction in incidence of type 2 diabetes with lifestyle intervention or metformin.
NEJM \textbf{346}(6), 393--403 (2002)

\bibitem{perkovic2019credence}
Perkovic, V., et al.:
Canagliflozin and renal outcomes in type 2 diabetes.
NEJM \textbf{380}(24), 2295--2306 (2019)

\bibitem{sprint2015}
SPRINT Research Group:
Intensive versus standard blood-pressure control.
NEJM \textbf{373}(22), 2103--2116 (2015)

\bibitem{naci2018comparative}
Naci, H., et al.:
How does exercise treatment compare with antihypertensive medications?
British Journal of Sports Medicine \textbf{53}(14), 859--869 (2019)

\bibitem{spirtes2000causation}
Spirtes, P., Glymour, C., Scheines, R.:
Causation, Prediction, and Search, 2nd edn.
MIT Press (2000)

\bibitem{vickers2006dca}
Vickers, A.J., Elkin, E.B.:
Decision curve analysis.
Medical Decision Making \textbf{26}(6), 565--574 (2006)

\bibitem{sritara2003egat}
Sritara, P., et al.:
Twelve-year changes in vascular risk factors: the EGAT Study.
Int.\ J.\ Epidemiology \textbf{32}(3), 461--468 (2003)

\bibitem{aekplakorn2011thai}
Aekplakorn, W., et al.:
Prevalence and management of diabetes in Thai adults: NHES IV.
Diabetes Care \textbf{34}(9), 1980--1985 (2011)

\end{thebibliography}

\end{document}
